% META: {"id": "TCOMPL-LOG-EX-026", "chapitre": "TCOMPL-LOGARITHME-HISTORIQUE", "type_objet": "exercice", "capacites": ["TCOMPL-LOG-C1"], "capacites_codes": ["C1"], "methodes": ["M3"], "parcours": 2, "competences": ["modeliser","raisonner"], "duree_min": 25, "mode_creation": "ex_nihilo", "sources_inspiration": [], "fichier_tex": "chapitres/TCOMPL-LOGARITHME-HISTORIQUE/exercices/TCOMPL-LOG-EX-026.tex", "corrige_tex": "chapitres/TCOMPL-LOGARITHME-HISTORIQUE/corriges/TCOMPL-LOG-CO-026.tex", "status": "approved"}
% BEGIN-VERIFY
% from sympy import *
% a, b, c = symbols('a b c', positive=True)
% # AM-GM pour 3 variables: (a+b+c)/3 >= (abc)^(1/3)
% # par concavite de ln
% lhs = ln((a+b+c)/3)
% rhs = (ln(a) + ln(b) + ln(c))/3
% # numerique
% assert lhs.subs({a:1, b:2, c:4}) == ln(Rational(7,3))
% assert simplify(rhs.subs({a:1, b:2, c:4}) - ln(2)) == 0
% assert ln(Rational(7,3)) > ln(2)
% END-VERIFY
\begin{exercice}{TCOMPL-LOG-EX-026}{3}{25}
Generaliser l'inegalite arithmetico-geometrique a trois variables positives :
\[ \frac{a+b+c}{3} \geq \sqrt[3]{abc}. \]
On demontrera le resultat en utilisant la concavite de $\ln$.

\end{exercice}
